% Figure: AM1.5G solar spectral irradiance versus wavelength, with the
% silicon band-gap cutoff at 1.12 eV (1107 nm) marked, separating the
% two single-junction loss regions: photons below the gap (unabsorbed,
% to the right) and the thermalisation loss of the excess energy of
% above-gap photons (to the left of the cutoff).
\begin{figure}[htbp]
\centering
\begin{tikzpicture}
\begin{axis}[
  width=12cm, height=7cm,
  xlabel={wavelength $\lambda$ (nm)},
  ylabel={spectral irradiance (W\,m$^{-2}$\,nm$^{-1}$)},
  xmin=300, xmax=2400, ymin=0, ymax=1.7,
  axis lines=left,
  tick label style={font=\scriptsize},
  label style={font=\small},
  clip=false,
]
% Schematic AM1.5G envelope: rises to a peak near 500 nm then falls with
% a long infrared tail. Sampled coordinates of the standard shape.
\addplot[engblue, very thick, smooth] coordinates {
  (300,0.05) (350,0.30) (400,0.65) (450,1.20) (500,1.55)
  (550,1.50) (600,1.45) (650,1.40) (700,1.30) (800,1.10)
  (900,0.95) (1000,0.78) (1107,0.62) (1200,0.55) (1400,0.30)
  (1600,0.28) (1800,0.18) (2000,0.12) (2200,0.08) (2400,0.05)
};
% silicon band-gap cutoff: 1.12 eV -> 1107 nm
\draw[engred, dashed, thick] (axis cs:1107,0) -- (axis cs:1107,1.55);
\node[engred, font=\scriptsize, anchor=south, align=center]
  at (axis cs:1107,1.55) {Si gap\\$1.12$ eV\\($1107$ nm)};
% loss-region labels
\node[enggray, font=\scriptsize, align=center] at (axis cs:600,0.45)
  {above gap:\\excess energy\\thermalises};
\node[enggray, font=\scriptsize, align=center] at (axis cs:1750,0.78)
  {below gap:\\not absorbed};
\end{axis}
\end{tikzpicture}
\caption[The AM1.5G solar spectrum and the silicon gap cutoff]{The standard
terrestrial solar spectrum (AM1.5G), spectral irradiance against wavelength,
with the silicon band-gap cutoff marked. A single-junction silicon cell loses
two ways at once. Photons to the right of the cutoff carry less than the gap
energy and pass through unabsorbed. Photons to the left carry more than the
gap energy, create one electron-hole pair each, and waste the excess as heat
as the carriers relax to the band edge. The two losses pull the optimum gap
in opposite directions and set the single-junction ceiling that
Figure~\ref{fig:vol04ch12:shockley-queisser} plots against gap.}
\label{fig:vol04ch12:am15-spectrum}
\end{figure}
